\scanpage{597}

\mspage{6}

\unsure{.} Reste à prouver que l'application (\unsure{a|A}) est \unsure{surjectif|surjective}. Soit donc $\unsure{\mathfrak m|m}$ un idéal maximal de $k_E$, et soit

\[
\Phi\overset{\mathrm{def}}{=}\{A\in\unsure{\mathcal P|\mathfrak P}(E)\mid \deleted{\unsure{f_A\notin\mathfrak m}}\ \text{\unsure{i.e.|t.q.} }\unsure{\chi|X}(\unsure{f_A|\varphi_A})\deleted{\unsure{\ne0}}\unsure{=1|\ne1}\},
\]

où $\chi:k_E\mathrel{\unsure{-|\longrightarrow}} k_E/\mathfrak m\inserted{=k(\mathfrak m)}$ est l'application canonique.

Je dis que $\Phi$ est un ultrafiltre. En effet, a) $\varnothing\notin\Phi$ car $f_\varnothing=0$ \inserted{d'où $\chi(f_\varnothing)=0\ne\deleted{\unsure{0}}1$ dans $k_\Phi$}, b) si $A,B\in\Phi$, on a \deleted{\unsure{$f$}} $A\cap B\in\Phi$, car $f_{A\cap B}=f_A\cdot f_B$ d'où $\chi(f_{A\cap B})=\chi(f_A)\chi(f_B)\deleted{\unsure{\ne}}\unsure{=1|\ne1}$,

c) Si $A\subset B$ et $A\in\Phi$, alors $B\in\Phi$, \deleted{en effet i.e. $B\notin\Phi\Rightarrow A\notin\Phi$} \deleted{on a $f_A f_B=\unsure{f_A}$,} \deleted{ou encore $\chi(f_B)=0\Rightarrow\chi(f_A)=0$.} En effet, on a $f_A=f_A f_B$, d'où $\chi(f_A)=\chi(f_A)\chi(f_B)$, \inserted{et comme $\chi(f_A)=1$, on trouve $\chi(f_B)=1$}

\deleted{\unsure{D'où} la \unsure{condition|conclusion}} d) On a $E\in\Phi$, i.e. $\chi(f_E)\deleted{\unsure{\ne}}=1$, or $f_E=1$ donc $\chi(f_E)=1$. \deleted{\unsure{Enfin},}

Ainsi, $\Phi$ est un filtre, \deleted{\unsure{puisqu'on a jusqu'ici} \unsure{utilisé} $\ne0$, et} \inserted{on n'a utilisé que le fait que $A/\mathfrak m$ est \deleted{\illegible{2}} un anneau $\ne0$ (pour \unsure{conclure} \unsure{que} $\varnothing\notin\Phi$).} \deleted{on \unsure{voit ici}} \deleted{\unsure{il faut|pour} un corps} \deleted{\unsure{Utilisons} Mais} \inserted{$A/\mathfrak m$ est intègre,} Prouvons que \deleted{si \unsure{$\mathfrak m$ est maximal},} \deleted{donc} $\Phi$ est un \uline{ultrafiltre}, ie. \unsure{a)} si $A\notin\Phi$, alors $B\overset{\mathrm{def}}{=}E\setminus A\in\Phi$.

En effet on a \deleted{\unsure{b)}} $f_A+f_B=1$, \deleted{$f_Af_B=0$, d'où $\chi(f_A)+\chi(f_B)=1$} \deleted{$\chi(f_A)\chi(f_B)=0$,} \inserted{$f_A f_B=0$, d'où \textcircled{a} $\chi(f_A)+\chi(f_B)=1$, et \textcircled{b} $\chi(f_A)\chi(f_B)=0$;} \deleted{Comme $\chi(f_A)=\unsure{0|2}$, c'est $\chi(f_B)=1$} $\chi(f_A)\ne1$, on conclut de \textcircled{a} $\chi(f_B)\ne0$, d'où de \textcircled{b} $\chi(f_A)=0$, d'où de \textcircled{a} $\chi(f_B)=1$. D'où $B\in\Phi$.

% archive sheet number: 596
